% SGA 3，第 VII_B 篇，英文重建稿——第 1 轮。
% 范围：权威底本分篇第 80 页，第 4.4.2 节全文，止于第 5 节之前。
% 合并阅读本页码： 593.
% 再版印刷页码： 582.
% 本范围内编者注： 141.
% 权威底本： Polo--Gille Exp7B-13oct24.pdf,
% SHA-256 D303690E0195654BBE7C3DAA5D89EA23E24CD709A3575217234AC19AE8836694.
% 曾参考 Jacob Reinhold 的 CC BY 4.0 Markdown 作为比对谱系。
% 原英文 Loop 1 阶段说明：Frobenius/高度归约图采用直接取自 PDF 的精确裁片。当前中文工程已采用原生图。

\subsubsection*{4.4.2}
\phantomsection
\label{sga3:VIIB:4.4.2}

继续假定 \(k\) 为特征 \(p\) 的域。设 \(G\) 为无穷小形式 \(k\)-群，仿射代数为 \(A\)。由于 \(G\) 是无穷小的，\(A\) 的每个开理想对充分大的 \(n\) 都包含
\[
  I_A^{\{p^n\}}
\]
因此，\(A\) 是群 \(\operatorname{Fr}_nG\) 的仿射代数
\[
  A/I_A^{\{p^n\}}
\]
的投射极限\textup{（参见 \SGARefLink{sga3:VIIA:4.1.2}{4.1.2}）}。等价地，每个无穷小形式 \(k\)-群都是有限高度形式 \(k\)-群的归纳极限。

假定 \(G\) 的高度至多为 \(n\)，记
\[
  H=G/\operatorname{Fr}G.
\]
\footnote{\label{sga3:VIIB:nde:141}编者在下文补充了原论证的细节。}
由 \SGARefLink{sga3:VIIB:2.4}{2.4} 和 \SGARefLink{sga3:VIIB:2.4.1}{2.4.1}，态射
\[
  \operatorname{Fr}:G\longrightarrow G^{(p)}
\]
分解为满态射
\[
  \pi:G\longrightarrow H
\]
后接单态射
\[
  i:H\longrightarrow G^{(p)}.
\]
于是有如下交换图：
\begin{center}
  \phantomsection
  \label{sga3:VIIB:4.4.2:height-reduction-frobenius-diagram}
  % SGA3-VII-DIAG-124
  % SGA3-VII-DIAG-124
% Authority: Exp7B-13oct24.pdf, local p. 80.
% Crop: figures/Exp7B_localp080_height_reduction_frobenius_diagram.png
% Locator: tex/components/73_expose_VIIB_section442_finite_height_filtration_en.tex, Loop 1 line 55.
\begin{tikzcd}[row sep=2.8em,column sep=3em]
G \arrow[r,"\pi"] \arrow[dr,"\operatorname{Fr}"']
  & H \arrow[r,"\operatorname{Fr}^{\,n-1}"] \arrow[d,"i"]
  & H^{(p^{\,n-1})}
      \arrow[d,"i^{(p^{\,n-1})}" right] \\
{} & G^{(p)} \arrow[r,"\operatorname{Fr}^{\,n-1}"]
  & G^{(p^n)}
\end{tikzcd}

\end{center}
由于函子 \(X\mapsto X^{(p)}\) 与纤维积交换，态射
\[
  i^{(p^{\,n-1})}
\]
仍为单态射。此外，\(\operatorname{Fr}^{n}(G/k)\) 为零，而 \(\pi\) 为满态射。因此，
\[
  \operatorname{Fr}^{n-1}\bigl(G^{(p)}/k\bigr)\circ i
  =
  i^{(p^{\,n-1})}\circ
  \operatorname{Fr}^{n-1}(H/k)
\]
为零。由于 \(i^{(p^{\,n-1})}\) 为单态射，
\[
  \operatorname{Fr}^{n-1}(H/k)=0,
\]
所以 \(H\) 的高度至多为 \(n-1\)。

由此可见，每个有限高度形式 \(k\)-群都有一个合成列，其商的高度至多为 \(1\)，因而可用 \(k\) 上的 \(p\)-Lie代数描述。

最后，由 \SGARefLink{sga3:VIIB:4.4.A}{注~4.4.A}，\(G\) 的仿射代数 \(A\) 是
\[
  k[[\omega_{G/k}]]
\]
的商。若 \(\omega_{G/k}\) 在 \(k\) 上有限维，则所有代数
\[
  A/I_A^{\{p^n\}}
\]
都是有限维 \(k\)-向量空间。因此，在特征 \(p>0\) 的域上，每个满足 \(\omega_{G/k}\) 在 \(k\) 上有限维的无穷小形式群 \(G\) 都是有限形式群的归纳极限\textup{（即有限长度的群；参见 \SGARefLink{sga3:VIIB:1.2.6}{1.2.6}）}。
